\documentclass[11pt,reqno]{amsart}
\usepackage[T1]{fontenc}
\usepackage{lmodern}
\usepackage{microtype}
\usepackage{mathtools}
\usepackage[letterpaper,textwidth=6.05in,textheight=8.62in,centering]{geometry}
\usepackage[hidelinks,pdfencoding=auto]{hyperref}
\usepackage[nameinlink,noabbrev]{cleveref}
\hypersetup{pdftitle={Unique tangent varifolds for planar network flow},
  pdfauthor={Kaiwen Jin},
  pdfsubject={Curvature flow of planar networks and tangent varifold uniqueness}}
\newtheorem{theorem}{Theorem}
\newtheorem{lemma}[theorem]{Lemma}
\ifdefined\pdftrailerid\pdftrailerid{}\fi
\numberwithin{equation}{section}
\allowdisplaybreaks[3]
\providecommand{\tightlist}{\setlength{\itemsep}{3pt}\setlength{\parskip}{0pt}}
\raggedbottom

\title[Unique tangent varifolds for planar network flow]{Unique tangent varifolds for planar network flow}
\author{Kaiwen Jin}
\thanks{\href{mailto:yui.ui.siki@gmail.com}{\texttt{yui.ui.siki@gmail.com}}.}
\keywords{Network flow, tangent varifold, curvature flow, triple junction}
\subjclass[2020]{53E10, 35K93, 49Q15}

\begin{document}
\begin{abstract}
We prove that a finite embedded planar network evolving by curvature has
a unique parabolically rescaled arclength varifold at each interior
finite-time singularity. The network has balanced triple junctions and
fixed boundary endpoints. A Gaussian length estimate, obtained by
coupling a finite geodesic description of the network inside a disk
to a signed estimate outside it, gives convergence of the spatial
weights. A local curvature-action estimate controls the projections of
the junctions swept out during short comparison windows. The scalar zero-number
theorem then gives enough intersections with transverse lines to
identify the tangent directions at arbitrary times. Both the
multiplicities and the directions are determined in the original
coordinates.
\end{abstract}
\maketitle

\section{Introduction}
\label{sec:introduction}

Planar network flow describes the motion of interfaces with triple
junctions, including grain boundaries in a two-dimensional idealization.
Parabolic rescaling at a singularity isolates its local geometry.
Different sequences of observation times could a priori give different
limits, even if the limiting lengths in each region were the same.
AIM Problem 159 \cite{AIM} asks for uniqueness as an arclength
varifold, retaining integer multiplicities and comparing limits in a
fixed frame.

The question is related to the blow-up problem of Mantegazza, Novaga,
Pluda and Schulze \cite[Open Problem~7.25 in the linked author version]{MNPS}.
Their problem also includes full-family local \(C^1\) convergence.
For a network without endpoints, Pluda and Pozzetta
\cite[Theorem~1.1]{PP} prove smooth convergence up to
reparametrization when there is a nonempty nondegenerate blow-up.
That blow-up is a compact regular shrinker on the same graph. Their proof uses a \L{}ojasiewicz--Simon inequality
in a smooth chart about that shrinker. At an arbitrary singularity,
some material edges may shrink to zero and exterior strands can
extend beyond every fixed disk in the rescaled variables. We use a
finite-dimensional energy that is subanalytic at collisions, together
with a separate estimate for the exterior strands. The conclusion
uses the local weak topology of arclength varifolds.

\begin{theorem}\label{thm:main}
Let \(\Gamma(t)\) be a finite embedded planar network moving by
curvature, with only \(120^\circ\) triple junctions and fixed endpoints
on the boundary of a smooth strictly convex domain. At any interior
finite-time singularity \((x_0,T)\), the arclength varifolds of
\begin{equation}\label{eq:rescale}
S_\tau=\frac{\Gamma(t)-x_0}{\sqrt{2(T-t)}},
\qquad \tau=-\tfrac12\log(T-t),
\end{equation}
converge locally weakly as \(t\uparrow T\) to a Gaussian-stationary
integral varifold carried by shrinker arcs. In particular, all locally
convergent subsequences have the same limiting arclength varifold,
including its integer multiplicities in the fixed ambient coordinates.
\end{theorem}

On a regular time interval, the network has fixed finite
incidence and compact material parameter intervals. Regular edges
have normal velocity equal to their curvature, the incident edges
have a common velocity at each junction, and the prescribed
\(120^\circ\) angles give zero sum of the three outward unit
tangents. The domain may be unbounded.

The arclength varifold of a network is the sum, over its material
edges, of arclength pushed forward by position and its unoriented
tangent line. Local weak convergence is tested on compactly supported
continuous functions on \(\mathbb R^2\times G(2,1)\), where \(G(2,1)\)
is the space of unoriented lines through the origin. Its spatial
weight is the marginal on \(\mathbb R^2\).
We use periodic parameters for a smooth closed component and mark a
point when an incidence description is needed.
We use smooth geometric representatives at positive times.
For each connected component other than a smooth closed curve, the regular-flow
conventions of \cite[Definitions~2.1, 2.6 and 2.11 and Theorem~5.6 of
the arXiv version]{MNPS} give this regularity. For a closed component,
apply Angenent's closed-curve regularity
\cite[Theorems~3.1 and 8.1 and their proofs]{AngCurve} with Euclidean
metric and speed \(V(\mathbf t,k)=k\). At an earlier regular time the
curvature of a compact regular loop is square integrable. Theorem~8.1
gives Hölder curvature at positive times, and uniqueness and
Theorem~3.1 give smooth geometric representatives after restarting.
Geometric reparametrization preserves the arclength varifold.

The proof has two parts. We first use planar face areas to rule out
collapsed cycles. The face potentials associated with Gaussian flux
then control the entire exterior of a fixed disk. The resulting
gradient estimate inside the disk is strong enough to give finite
rescaled action and convergence of the spatial weight. We next compare
an arbitrary starting slice to a nearby small-defect slice. On a
regular arc of the limiting support, an estimate for the projection
of the swept junctions allows us to use transverse intersection
counts to identify the limiting tangent lines. This final step makes
the conclusion a varifold statement.

\section{Gaussian estimates on regular slices}\label{sec:gaussian}

\subsection{Weighted variation and local length ratios}

In rescaled variables write \(y\in S_\tau\), \(r=|y|\), \(\rho(r)=e^{-r^2/2}\), and on each oriented edge use the unit tangent \(u\), normal \(n=Ju\) obtained by counterclockwise rotation by \(90^\circ\), and signed curvature \(k\). Let \[
 a=y\cdot n,\quad b=y\cdot u,\quad h=k+a,\quad f=a\rho,\qquad
 F=\int\rho\,ds,\quad D=\int h^2\rho\,ds .
\] All integrals without a specified domain are over the complete material network, once per edge. For velocity \(K n+\lambda u\), differentiating the speed in a fixed material parameter gives the length derivative \((\partial_s\lambda-kK)ds\). Consequently a time-dependent scalar weight \(\varphi\) gives interior integrand \(\partial_t\varphi+K n\cdot\nabla\varphi-kK\varphi\) and the oriented endpoint bracket of \(\lambda\varphi\). In the rescaled variables \(dt/d\tau=2(T-t)\), so the normal speed is \(h\). Since \(\nabla\rho=-y\rho\), this gives \[
 F'=-D+B_0,\qquad
 |B_0|\le B=\sum_{\text{fixed endpoints }p}|p_\tau|\rho(|p_\tau|),\qquad
 p_\tau=\frac{e^\tau(p-x_0)}{\sqrt2}.
\] At a true vertex the outward tangent sum vanishes and the velocity is common, so the entire boundary contribution cancels. At a physical endpoint \(p\), \(B_0\) uses the term \(\rho(|p_\tau|)p_\tau\cdot u_{\rm end}\), with \(u_{\rm end}\) pointing along the edge toward its endpoint. Regularity on the positive-time slices makes \(B_0\) continuous in time. Since \(x_0\) is interior, \(B\) decays superexponentially. Consequently \(F\to c\) and \(\int^\infty D<\infty\). Put \(g_0(\tau)=\int_\tau^\infty|B_0(s)|\,ds\) and \(H=F-c+g_0\). Then \(H\downarrow0\) and \(H'\le-D\). Write \(\mu_\tau\) for the spatial weight of the arclength varifold of \(S_\tau\). For each fixed \(\varphi\in C_c^1(\mathbb R^2)\), the fixed endpoints are outside its support at late times and the same calculation gives \begin{equation}\label{eq:weight-tests}
 \frac{d}{d\tau}\mu_\tau(\varphi)
  =\int\big[h n\cdot(\nabla\varphi+\varphi y)-\varphi h^2\big]\,ds,
 \qquad
 \big|\tfrac{d}{d\tau}\mu_\tau(\varphi)\big|
 \le C_\varphi(\sqrt D+D).
\end{equation}

For the centered length bound, fix a physical center \(z_0\) near \(x_0\), radius \(d_0>0\), and kernel time \(\theta=t+d_0^2\) at the physical slice \(t\) of interest. At a physical time \(v\le t\) put \(q=\theta-v\), \(Y=X-z_0\), \(a_0=Y\cdot n\), \(b_0=Y\cdot u\) and \(\Phi=(4\pi q)^{-1/2}e^{-|Y|^2/(4q)}\). The weighted length derivative above, with \(K=k\), has interior integrand \[
\left(\frac1{2q}-\frac{|Y|^2}{4q^2}-\frac{a_0k}{2q}-k^2\right)\Phi
=-\left(k+\frac{a_0}{2q}\right)^2\Phi
+\partial_s\!\left(\frac{b_0\Phi}{2q}\right),
\] using \(b_{0,s}=1+a_0k\) and \(\Phi_s=-b_0\Phi/(2q)\). The endpoint bracket of \(\Phi(\lambda+b_0/(2q))\) cancels at triple junctions and at a smooth mark. At a fixed endpoint \(\lambda=0\) and the total absolute contribution is bounded by \(\int_0^\infty C l q^{-3/2}e^{-l^2/(4q)}\,dq\), with \(l=|p-z_0|>0\), uniformly bounded by scaling \(q\) by \(l^2\). The calculation uses \(q>0\) only on the actual integration interval, even if \(\theta>T\). Integrate from a fixed regular time earlier than all the times \(t\) in question. The initial kernel integral is bounded by its finite physical length times a constant since \(q\) there has a fixed positive lower bound. At \(t\) the lower kernel bound is a positive numerical constant times \(d_0^{-1}\) on \(B_{d_0}(z_0)\), so its length is at most \(C d_0\), uniformly in \(d_0\) and in a fixed interior neighborhood of \(x_0\). Rescaling gives \(\mu_\tau(B_r(z))\le Cr\) for all \(r>0\) and all \(z\) in each fixed compact set for sufficiently late \(\tau\). Dyadic annuli give uniform Gaussian moments of every fixed polynomial degree. The bound in \eqref{eq:weight-tests} now follows by Cauchy--Schwarz on the compact test support.

\subsection{Face areas and Gaussian flux}\label{sec:faces}

Along any material edge \(a_s=-bk\), \(\rho_s=-b\rho\), and hence \(f_s=-b h\rho\). The length bound and Cauchy--Schwarz show that the oscillation of \(f\) on an edge is \(O(\sqrt D)\). Orient the material edges once. Representative fluxes along the edges have incidence residual \(O(\sqrt D+B)\), by the outward tangent sum at each triple vertex and the bound at a fixed endpoint. Project the representative vector to the real kernel of the finite incidence matrix, choosing a fixed bounded linear projection, and denote the resulting circulation by \(c_e\) on edge \(e\). Then \[
 \sup_{y\in e}|f(y)-c_e|\le C(\sqrt D+B).
\] A marking of an isolated loop adds a bivalent point with no incidence residual. In a finite embedded planar graph the oriented boundaries of the bounded faces are a basis of the circulation space. To see spanning, decompose a circulation into directed cycles by orienting edges with positive flux and successively subtracting the minimum positive flux on a cycle. Each simple oriented cycle is the signed sum of bounded face boundaries enclosed by it. Independence follows by connecting each face to the unbounded face in the dual and testing the differences on all edges. Thus \(c_e=\pi_{\rm left}-\pi_{\rm right}\) for potentials \(\pi_F\) with exterior potential zero.

For a bounded complementary face \(F_1\) use all oriented boundary walks with the face on the left. The sum \(K_{F_1}\) of their smooth signed curvatures belongs to \((\pi/3)\mathbb Z\) and is constant throughout the regular evolution. Follow the oriented boundary walks by taking the next edge on the face side at each vertex. Lift the tangent angle on each edge and join its end tangents by the signed turn at the corner through the face angle. The total over a closed walk is an integer multiple of \(2\pi\), since the final tangent agrees with the initial tangent. At each \(120^\circ\) or \(240^\circ\) sector the added turn is \(+\pi/3\) or \(-\pi/3\). A bridge used twice gives opposite contributions of its smooth curvature. If the same face uses several sectors at a vertex, the oriented corners are counted once for each sector passage. If a boundary walk uses a degree-one endpoint, it goes out and back on the same edge and turns by \(\pi\) modulo \(2\pi\). This contributes a multiple of \(\pi/3\) to the turn calculation. The two smooth curvature integrals cancel. At a fixed endpoint the area-variation bracket is zero as well. The physical area is \(\frac12\sum\int\det(X,X_v)\,dv\) on the same oriented walks by Green's formula, using the material parameter \(v\), including hole orientations and cancelling bridge integrals. Differentiation and integration by parts give \(\int\det(X_t,X_v)\,dv+\frac12[\det(X,X_t)]\) on each piece. The boundary brackets cancel at successive joins because position and velocity agree. Hence the area derivative is \(-K_{F_1}\), using \(\det(n,u)=-1\) for the inward normal. In particular \begin{equation}\label{eq:face-area}
 A_{F_1}^\mathrm{resc}(\tau)
 =\tfrac12 e^{2\tau}A_{F_1}(T)+\tfrac12 K_{F_1}.
\end{equation} A face with \(A_{F_1}(T)=0\) has positive rescaled area at least \(\pi/6\). Let \(A_{\rm col}\) be the sum of their rescaled areas. A face with \(A_{F_1}(T)>0\) has rescaled area growing as \(e^{2\tau}\) and reaches radius \(l_{F_1}=a_{F_1}e^\tau\) for some \(a_{F_1}>0\). To estimate its potential, start in that face at radius \(>l_{F_1}\), take a path to infinity inside the open exterior of that disk, and choose transverse crossings avoiding the true vertices. Loop erasure in the finite dual bounds the number of retained crossings. Each has \(|f|\le l_{F_1}\rho(l_{F_1})\) for late times. This yields \[
 |\pi_{F_1}|\le C(\sqrt D+B+\beta),\qquad
 \beta=\sum_{F_1:A_{F_1}(T)>0} l_{F_1}\rho(l_{F_1})
\] for every such persistent face. At all sufficiently late times every bounded face has rescaled area at least \(\pi/6\). An embedded simple material cycle encloses at least one bounded face, so its length has a uniform positive lower bound. A loop of length \(\ell\) and its bounded interior lie in the closed disk of radius \(\ell\) centered at any point of the loop. The area of this disk is \(\pi\ell^2\). This bound does not bound the length of an individual edge from below.

It will be convenient to use the positive quantity \(d=\sqrt D+\eta\), where \(\eta=B+\beta+\sqrt{B+g_0}+e^{-\tau}\) is integrable at late times. The last term keeps \(d\) positive even at a zero-defect slice of a network without endpoints.

\subsection{Small-defect limits}\label{sec:compactness}

Choose a sequence \(\tau_j\to\infty\) with \(D(\tau_j)\to0\). The local length ratios and \(|k|^2\le 2h^2+2|y|^2\) give local \(L^2\) curvature bounds on these slices. In each larger ball, retain the components of the edge portions in that ball that meet a strictly smaller concentric ball. Their number is uniformly bounded. A piece is a whole material edge, or it traverses the collar with length at least its thickness. Arclength and Cauchy--Schwarz give \(C^1\) and weak \(W^{2,2}\) convergence on the pieces of positive limiting length, after extraction. Pieces of vanishing length give no mass. Exhausting the plane gives a local integer arclength varifold \(V_*\) with spatial weight \(\mu_*\). Each surviving arc solves \(h=0\) by weak convergence of curvature and strong convergence of position and tangent, so it is analytic by its ordinary differential equation. The full limit is Gaussian stationary as well. Against a compactly supported smooth vector field the weighted first variation at \(\tau_j\) is the integral of \(-h n\) against that field with weight \(\rho\). The vertex terms vanish, and the absolute value tends to zero by \(D(\tau_j)\to0\). The first variation integrand is continuous in position and tangent line, so it passes to the varifold limit. The same compactness applies on every sequence with \(D\to0\). On each nonvanishing piece the tangent modulus is \(|u(s)-u(s')|\le |s-s'|^{1/2}(\int k^2\,ds)^{1/2}\). Use Arzelà--Ascoli on the arclength positions and tangents, parametrizing on common subintervals and allowing the endpoints to converge. The shrinker equation follows by integrating \(u_s=-(y\cdot n)n+hn\), using local \(L^1\) convergence of \(h\) to zero. Its analytic ODE extends the surviving arcs past their endpoints. In each smaller compact ball the finite analytic arcs meet at isolated transverse crossings or agree as geometric germs when their tangent lines agree, by ODE uniqueness and reversibility. Along with the finite endpoint set this gives a regular arc patch of constant positive integer weight almost everywhere. A coincident open arc contributes once per material arclength piece.

\section{The Gaussian length and force outside a disk}\label{sec:exterior}

\subsection{Gaussian length of the exterior}

Along the chosen sequence extract limits of the finitely many \(c_e\). If the limit on an edge is nonzero, that edge lies in a fixed ball because \(|f|\le r\rho(r)\) and the error to \(c_e\) vanishes. On an edge with zero limiting flux, every surviving arc has \(a=0\) and is radial wherever \(y\ne0\). Therefore \(\operatorname{spt}\mu_*\) consists locally of radial segments at radii greater than some \(R_0\). In particular, it contains no nonradial regular arc there.

Fix \[
 R_0+2<R_-<R,\qquad R\ge4,\qquad \pi R^2>A_{\rm col}.
\] Take \(P\) to be the maximum of \(|\pi_F|\) over bounded faces meeting \(\{r>R_-\}\), with maximum zero when there are none. Every edge point in this open exterior has adjacent faces meeting it. Thus throughout this exterior, in particular on a neighborhood of the circle at \(R\), \begin{equation}\label{eq:flux}
 |f|\le 2P+Cd.
\end{equation}

For any sequence of late slices with \(D\to0\) and \(\mu_\tau\to\mu_*\), we have \(P\to0\). Otherwise use an exterior dual path from a face whose potential is bounded away from zero, as above. Some crossed edge has limiting nonzero circulation, along a subsequence with fixed label. The whole edge is in a fixed bounded ball. If its length does not tend to zero, a fixed short interval at the crossing on one side, within \(r>R_0\), yields a nonradial regular limit arc by the small-defect compactness. If the crossed edge collapses, both endpoints converge to the same point \(y_0\) with \(|y_0|\ge R_-\). The cycle decomposition supplies a simple directed cycle through this edge, with nonzero limiting circulation on every edge. If the next edge of the cycle kept positive length, a short initial arc based at the shared endpoint would lie in \(r>R_0\) and yield a nonradial limit. Thus the next edge also collapses to \(y_0\). Repeating the implication at the finitely many successive endpoints makes the whole cycle collapse there, contrary to the area lower bound from \eqref{eq:face-area}. This proves the claim. Exterior dual paths here can be chosen in the open set \(\{r>R_-\}\), whose closure at any retained crossing is separated from \(R_0\).

Take \(\Delta=C_1(P+d)\) with \(C_1\) fixed sufficiently large and let \(S\) be the large root of \(S\rho(S)=K\Delta\) for one sufficiently large fixed \(K\). The following estimates concern \(P+d\to0\), so \(S\to\infty\) and \(S^2\asymp\log(1/(P+d))\). In the annulus \(R\le r\le S\), \eqref{eq:flux} yields \(|a|/r\le 1/4\) by taking \(K\) large. In particular \(|b|/r\ge 1/2\), and a true vertex is excluded there because its three directions cannot all be within these two opposite radial caps. No fixed endpoint lies in the annulus because its rescaled \(r\rho(r)\) would contribute at least \(K\Delta\) to \(B\le d\). The portions in the closed annulus are radial monotone through-arcs, and there are \(N\) crossings at both spheres with \(N\le C\), by the length ratio on \(B_S\) and the length of at least \(S/2\) in \([\tfrac12 S,S]\) on each through-arc. The collar in \eqref{eq:flux} excludes true vertices or tangencies exactly at the clipping circle as well.

Define \(G(r)=\int_r^\infty\rho(v)\,dv\), \(z(r)=G(r)/(r\rho(r))\), and \(w(r)=-z'(r)/r\) for \(r\ge R\). The representation \(G/\rho=\int_0^\infty e^{-ru-u^2/2}\,du\) gives \(w\asymp r^{-4}\) and \(z^2\le Cw\) by differentiation. On the annulus \(|f|\le C\Delta\) and \(ds/dr\le2\), so \[
 \int_{R\le r\le S} w a^2\rho\,ds
 \le C\Delta^2\int_R^S r^{-4}e^{r^2/2}\,dr
 \le C\Delta^2 S^{-5}e^{S^2/2}
 \le C\Delta S^{-4}.
\] The middle estimate follows by differentiating \(r^{-5}e^{r^2/2}\) with \(R\ge4\). On the complete exterior at \(S\), integrate \((b\rho)_s=(1-r^2+ah)\rho\), including all incident legs at each exterior true vertex. The signed endpoint sums at the transverse cuts are bounded by \(NS\rho(S)\), and the fixed endpoint term is bounded by \(B\). Since \(|a|\le r\), Young's inequality absorbs a fixed fraction of \(\int_{r\ge S} r^2\rho\,ds\) in the term with \(ah\), giving \[
 \int_{r\ge S}r^2\rho\,ds\le C(NS\rho(S)+D+B)\le C\Delta .
\] It follows that \(\int_{r\ge S}w a^2\rho\,ds\le C\Delta S^{-4}\) as well. These integrations include any exterior component with no cut.

From \(r z'=r^2z-z-1\) we obtain \((zb\rho)_s=-\rho+w a^2\rho+z h a\rho\). Integration on the complete exterior of \(R\) now gives \[
 \int_{r\ge R}\rho\,ds
 =G(R)\sum_{r=R}\sqrt{1-(a/R)^2}
  +\int_{r\ge R}w a^2\rho\,ds+\int_{r\ge R}z h a\rho\,ds+E_{\rm end},
\] with \(|E_{\rm end}|\le\sum_p G(|p_\tau|)\le C B\) for late times. On \(r=R\), \(a=O_R(\Delta)\) by \eqref{eq:flux}, so the cut sum differs from \(NG(R)\) by \(O_R(\Delta^2)\). The nonnegative \(wa^2\) integral is \(O(\Delta S^{-4})\). The curvature cross term has absolute value at most
\[
C\sqrt{D}\left(\int_{r\ge R} w a^2\rho\,ds\right)^{1/2}
 \le C d\sqrt{\Delta}\,S^{-2}
 \le C\Delta S^{-4}.
\]
The last inequality follows from \(d\le C\Delta\) and \(\sqrt{\Delta}S^2\to0\), which also gives \(\Delta^2=O(\Delta S^{-4})\). Together with the endpoint bound these estimates control both signs of the length difference. \begin{equation}\label{eq:outer-length}
 \left|\int_{r\ge R}\rho\,ds-N G(R)\right|
 \le C\left[B+\frac{P+d}{\log^2(1/(P+d))}\right].
\end{equation}

\subsection{Face layers and the angular boundary variation}\label{sec:force}

In this section use the regime \(P/d\) larger than a sufficiently large fixed constant. When some faces are selected, keep as boundary edges exactly those with one selected and one unselected incident side. Orient them with the selected side on the left. A bridge has the same face on both sides and contributes no boundary edge. At a true vertex, the cyclic list of three face sectors changes selected status zero or two times. The boundary edges thus form disjoint embedded simple closed curves. The exterior face is unselected.

All positive face potentials are \(O(d)\). If the positive maximum is much larger, choose a level between half the maximum and the maximum, at distance much greater than \(d\) from every face potential. Select the faces whose potentials exceed that level. Along their oriented boundary the actual \(f\) is strictly positive by the circulation approximation. An outer boundary is counterclockwise. On any such cycle the polar angle satisfies \(\phi_s=-f/(r^2\rho)<0\), contradicting its winding about the origin. No such boundary contains the origin on an edge or vertex, since \(f\) there would vanish. This rules out the positive maximum in that regime.

For a negative threshold \(s>0\), select all faces with \(\pi_F<-s\) using the boundary convention just given. Restrict to \(0<s<P\), omitting \(s\le C_0 d\) and intervals of radius \(C_0 d\) about all numbers \(-\pi_F\), for a fixed \(C_0\) large enough. The omitted measure is \(O(d)\). Across each retained interval \(I\) the selected union is fixed, all its faces collapse by the persistent-face bound, and actual oriented flux is strictly negative on its boundary. Thus the polar angle strictly increases along each oriented boundary with its continuous lift at the corners. Every boundary has positive winding and contains the origin in its bounded interior. A hole would have clockwise boundary, and disjoint components without holes cannot both contain the origin. A nonempty selection therefore has a single counterclockwise boundary. The angle increases strictly and its winding number is one, so the boundary meets each ray from the origin once. For context, Chang and Lue prove star-shapedness for the union of bounded regions of an exact regular shrinker \cite[Lemma~3.1]{CL}. Here the sign of the approximate flux gives star-shapedness for the selected region in the evolving slice. Its area is at most \(A_{\rm col}\).

Label the sphere cuts by local angle coordinates \(\alpha_a\). On the selected boundary, pair each outward crossing \(a\) with the next inward crossing \(b\), using the intervening exterior piece of that oriented boundary, including any true corners. Its positive angular increment is \[
 q_{I,ab}=\alpha_b-\alpha_a+2\pi n_{I,ab}>0.
\] There are finitely many possible pairings, signs, and winding integers in fixed angle charts. Indeed the number of retained intervals is bounded by the number of original faces, the cut count is bounded, and each piece's angular increment is at most \(2\pi\). Every such exterior piece has only its first and last leg in \(R\le r\le S\): transversality makes each annular leg monotone from one sphere to the other. Its angle increment there is at most \(C(P+d)\int_R^S dr/(r^2\rho(r))\le C S^{-2}\), using the integrand estimate as in the exterior length estimate with \(r^{-3}e^{r^2/2}\) in the derivative. The increment on \(r\ge S\) is at most \(2 A_{\rm col}/S^2\), by the polar area formula \(\frac12\int r(\phi)^2\,d\phi\) for the selected region. Thus every component of \(q\) is \(O(1/\log(1/P))\).

The face of magnitude \(P\) meeting \(r>R_-\) has negative potential in this regime. There is a subinterval \(J\) of \([P/2,P]\), of length at least \(c_G P\), at distance at least \(c_G P\) from every face level. Indeed, divide this interval by the finitely many face levels and take a middle subinterval of a longest piece. Shrinking the graph-dependent constant \(c_G>0\) gives both bounds. For sufficiently large \(P/d\), the interval \(J\) lies in a retained interval \(I\). The layer on \(I\) is constant, reaches \(r>R_-\), and has coefficient \(|I|/P\ge c_G\) in the angular pairing. If its maximal radius were \(\le R\), it would attain that maximum on its boundary at \(r\in(R_-,R]\). At a smooth radius maximum \(b=0\), hence \(|f|=r\rho(r)\). At a true vertex maximum \(r\rho(r)\le C\max |f|\) over its two nonparallel boundary-arm normals. Both contradict \eqref{eq:flux} as \(P+d\to0\), since \(r\rho(r)\) is bounded below on \([R_-,R]\). Thus this region goes beyond \(R\). It also has boundary inside \(R\), since it is star-shaped about the origin and has area at most \(A_{\rm col}<\pi R^2\). The collar lets us infer an actual crossing at \(R\) from a face reaching past \(R_-\).

At any level in \(J\) the adjacent potentials on each selected boundary edge lie on opposite sides of the level and are separated from it by at least \(c_G P\). Their difference therefore has magnitude at least \(2c_GP\). The flux approximation in \cref{sec:faces} then gives at least \(c'_GP\) for the magnitude of the actual oriented boundary flux, once the fixed threshold for \(P/d\) is large enough. Its magnitude is at most \(CP\) in the exterior by \eqref{eq:flux}. After an outward cut at \(R\), follow the exterior boundary up to its first smooth radial critical point or its first true vertex. One of these events must occur before the return. The same smooth or two-normal calculation implies \(r\rho(r)\le K_1P\) at the first event, so the outgoing radial-monotone arc reaches the large root \(S_1\) of \(S_1\rho(S_1)=K_1P\). On \([S_1-S_1^{-1},S_1]\), integrate \(\phi_s=-f/(r^2\rho)\) using \(ds/dr\ge1\). The angular increment is at least \(c P e^{S_1^2/2}/S_1^3\ge c' S_1^{-2}\). In particular, with \(\zeta=\|q\|\) and \(\lambda_{I,ab}=|I|/P\) for every pair in each retained interval, we have \begin{equation}\label{eq:pairings}
 \zeta\asymp\frac1{\log(1/P)},\qquad
 \|\lambda\|\le C,\qquad \lambda\cdot q\ge c_1\zeta .
\end{equation}

At each cut orient the core tangent out of the compact core. Its true Gaussian angle covector is \(\rho u\cdot\partial_\alpha y=-a\rho=-f\). Replace \(f\) by the circulation with error \(O(d)\). The negative potential at each adjacent face is the negative integral of its sublevel indicators in \(s\), and the positive potential is \(O(d)\). Every adjacent face at a cut belongs to the exterior face set unless it is the unbounded face, whose potential is zero. Hence layers at \(s>P\) contribute nothing at any cut. Layers omitted near face levels and zero also contribute \(O(d)\). At the outward start \(a\), the boundary tangent agrees with the core's outward-pointing tangent. The selected negative potential is on the left, so its contribution to \(-f\) is \(+1\) per unit level. At the return \(b\), these two tangent orientations are opposite, so the contribution is \(-1\). Their pair of coefficients is \(-\mathrm D_\alpha q_{I,ab}\). Thus the vector of core Gaussian angle covectors is \begin{equation}\label{eq:force}
       -P(\mathrm D_\alpha q)^T\lambda+O(d).
\end{equation} This calculation uses all original faces and all the actual exterior pairings. In the regime \(P=O(d)\), \eqref{eq:flux} directly bounds all angle covectors by \(O(d)\).

\section{Convergence of the spatial weights}\label{sec:weights}

\subsection{The finite-dimensional length}\label{sec:core}

All physical endpoints lie outside \(B_R\) at late times, since each has \(|p_\tau|=e^\tau|p-x_0|/\sqrt2\) with \(|p-x_0|>0\). Along any sequence approaching \(\mu_*\) with \(D\to0\), one may pass to a fixed clipped graph with a fixed number of division points on each material piece inside \(B_R\). There are finitely many possibilities because each piece is either a whole material edge or ends at a transverse cut, and the cut count is bounded. The length bound allows an equal-arclength subdivision of every piece into one common fixed number of subpieces with length below a prescribed short normal-coordinate radius. Use the coordinates \(x\) of the true vertices, division points and clipping angles, using fixed angular chart lifts along a subsequence.

On each circle chart the position is analytic in the angle. Treat loops by a mark and at least two divisions. The two endpoints of each positive-length subpiece are distinct, while subpieces sharing a vertex use the same coordinate there. All the coordinates remain in a bounded set and have a common subsequential limit. The functions \(q_{I,ab}\) are affine in these coordinates for a fixed pairing pattern.

In the conformal metric \(\rho^2|dy|^2\), replace each subpiece by the short geodesic joining its endpoints, and call \(E(x)\) the sum of those short distances plus \(NG(R)\). Distances here are the local distances in normal neighborhoods of a fixed compact set, with a margin inside the short normal-coordinate radius. The construction stays in the local normal balls of that compact set. The geodesic equation is an analytic first-order system in position and velocity and its exponential map is analytic in the initial data. For this analytic dependence, the Picard proof of \cite[Theorem~4.2]{Teschl} applies to the vector system: complexify on a smaller polydisc and use the bounded Jacobian in a product norm. On a sufficiently small common time disc the Picard iterates remain in the polydisc and converge uniformly on smaller polydiscs of time and initial data. Their holomorphic limit gives the analytic dependence. The derivative of the exponential map in the initial vector at zero is the identity. The analytic inverse function theorem gives the analytic log map near the diagonal. For \(Z(r,\omega)=\exp_p(r\omega)\) and \(|\omega|_g=1\), \(Z_r\) is a parallel unit tangent. Because the coordinate derivatives commute and the metric connection has zero torsion, \(\partial_r\langle Z_r,Z_\omega\rangle_g=\frac12\partial_\omega\langle Z_r,Z_r\rangle_g=0\). The inner product tends to zero at the center. Thus the radial component of any curve's velocity has magnitude at most its metric speed. The outward radial segment minimizes length inside the normal ball. Leaving the ball costs at least its radial radius. In particular the squared short distance is the squared metric norm of the log vector. The same radial calculation identifies its distance gradient away from the center. For the length differential of the radial geodesic, differentiate its length through its smooth endpoint family and integrate the covariant derivative of the variation along it. The interior term vanishes because the unit tangent is parallel, so the two endpoint covectors are the terminal tangent and the negative initial tangent in the metric.

Compactness permits the normal radii and the endpoint comparison constants to be restricted uniformly. Its square root is Lipschitz as a distance, and is globally subanalytic after restriction to a bounded chart with compact closure inside a slightly larger such chart. The resulting \(E\) is smooth near every observation with positive subpiece lengths. Coincident endpoints in limiting coordinates are allowed.

The length and covector errors are uniform even for subpieces tending to zero. To check this, let \(\ell\) be one subpiece's Gaussian arclength and \(\kappa_g=h/\rho\) its signed Gaussian geodesic curvature. The connection for \(g=e^{2\phi}|dy|^2\), \(\phi=\log\rho\), is \(\nabla^g_U V=\nabla_U V+U(\phi)V+V(\phi)U-(U\cdot V)\nabla\phi\) by the coordinate metric formula. Pair this with \(n/\rho\) in \(g\) for \(U=V=u/\rho\) to obtain \(\kappa_g=\rho^{-1}(k-\partial_n\phi)=h/\rho\). Compare its Gaussian arclength position and tangent up to length \(s\le\ell\) with the geodesic of the same initial data. In the smooth geodesic system on a compact chart the curve has an additional acceleration of size \(|\kappa_g|\), so Grönwall and Cauchy--Schwarz bound the tangent error by \(C\sqrt s\,\|\kappa_g\|_{L^2[0,s]}\) and the position error by \(C s^{3/2}\|\kappa_g\|_{L^2[0,s]}\). The second bound follows by integrating the first tangent difference. The log map at the initial point has bounded differential and the reference log vector has length \(s\). Thus the initial direction of the geodesic to the actual endpoint differs from the actual initial direction by \(O(\sqrt s\,\|\kappa_g\|_2)\). The length \(s_g\) of the log vector satisfies \(|s_g-s|\le Cs^{3/2}\|\kappa_g\|_{L^2[0,s]}\), so it is comparable to \(s\) for small defect. Comparing terminal tangents uses both the initial-direction error and this difference of evaluation lengths. Smooth dependence gives a tangent error of \(O(\sqrt s\,\|\kappa_g\|_{L^2[0,s]})\). The radial distance derivative along the actual subpiece is the cosine of its angle with the outward radial geodesic. Integrating the bound on \(1-\cos\) gives the Gaussian length excess \(O(\ell^2\|\kappa_g\|_2^2)\). On the core, \(\int\kappa_g^2\,ds_g=\int h^2/\rho\,ds\le C_R D\), where \(ds_g=\rho\,ds\).

The endpoint derivatives of the geodesic lengths are their Gaussian tangent covectors with outward signs. The actual covectors cancel at each artificial division and each true vertex. At the clipping angles, the actual ones are those of \eqref{eq:force}. Together with \eqref{eq:outer-length} this proves \begin{equation}\label{eq:core-length}
 \begin{split}
 |F-E(x)|&\le C\left[D+B+\frac{P+d}{\log^2(1/(P+d))}\right],\\
 \|\nabla E(x)+P(\mathrm Dq)^T\lambda\|&\le C d\qquad(P/d\text{ above the fixed threshold}),
 \end{split}
\end{equation} and \(\|\nabla E(x)\|\le Cd\) when \(P=O(d)\). The right side includes the exterior contribution through \eqref{eq:outer-length} and \eqref{eq:force}, with the boundary estimates valid on all exterior branches.

\subsection{Separation of the perturbed gradient}\label{sec:separation}

\begin{lemma}\label{lem:separation}
Let \(E\) be a globally subanalytic function on a bounded globally subanalytic coordinate domain, and let \(U_E\) be a globally subanalytic open subset where \(E\) is smooth. Fix an affine map \(q\), a value \(c\), and positive constants \(C_*,c_*,C^*\). Consider all \(x\in U_E\) and \(p,\lambda\) with
\[
 \zeta=\|q(x)\|\in(0,1),\qquad 0<p\le\zeta^M,\qquad
 \|\lambda\|\le C^*,\quad \lambda\cdot q(x)\ge c_*\zeta,\qquad
 |E(x)-c|^\theta\le C_*p
\]
for a rational \(0<\theta<1\) and an integer \(M>\theta/(1-\theta)\). Then for all sufficiently small \(\zeta\) there are \(c_2>0\) and \(0<\varepsilon\le 1/2\) such that
\begin{equation}\label{eq:separation}
 \|\nabla E(x)+p(\mathrm Dq)^T\lambda\|
 \ge c_2\,p\,\zeta^{1-\varepsilon}.
\end{equation}

\end{lemma}

\begin{proof}
We use real quantification, definable choice, and one-variable Puiseux expansions in the globally subanalytic structure \cite[Theorem~2.1.4 and Propositions~2.2.1 and 1.8.4]{Valette}. On each nonempty fiber at fixed \(\zeta\), let \(R(\zeta)\) be the infimum of \(\|\nabla E+p(\mathrm Dq)^T\lambda\|/(p\zeta)\) over the displayed constraints. The infimum uses the whole definable set of these coordinates and parameters. Both the domain of this infimum and its graph are definable in the globally subanalytic structure, by definability of the smooth gradient, the norm, rational powers, and quantification over definable sets. If there were arbitrarily small \(\zeta\) with bounded ratio, definable choice on the inequality with a fixed ratio bound would provide bounded curves \(x(\zeta),p(\zeta),\lambda(\zeta)\) on an interval down to zero. The convergent Puiseux expansions of these curves give \(\zeta x'(\zeta)=o(1)\), \(p(\zeta)\sim b_0\zeta^m\) with \(b_0>0\), \(m\ge M\), and \(q(\zeta)=\zeta q_0+O(\zeta^{1+\gamma})\) for some \(\|q_0\|=1\) and \(\gamma>0\), since \(\|q(\zeta)\|=\zeta\). Differentiating the convergent Puiseux expansion yields \(q'(\zeta)=q(\zeta)/\zeta+o(1)\). At each point in the chosen smooth observation set the chain rule gives, for sufficiently small \(\zeta\),
\[
 (E(x(\zeta)))'=-p(\zeta)\lambda(\zeta)\cdot q'(\zeta)+o(p(\zeta))
 \le -\tfrac12 c_*p(\zeta).
\]
Since \(E\to c\), integrate from zero to obtain \(|E-c|\ge c_3 p(\zeta)\zeta\) by the Puiseux order of \(p\). But
\[
 \frac{(p\zeta)^\theta}{p}
 \ge \zeta^{\theta-M(1-\theta)}\longrightarrow\infty,
\]
contradicting the power bound. The infimum ratio therefore tends to infinity on its nonempty fibers. If these fibers persist arbitrarily near zero, their definable domain contains an initial interval. The Puiseux expansion of the reciprocal of the infimum on that interval gives \(R(\zeta)\ge c_2\zeta^{-\varepsilon}\) for some \(\varepsilon>0\), which proves \eqref{eq:separation}.
\end{proof}

\subsection{Trapping and spatial convergence}\label{sec:trapping}

Consider any sequence of late slices converging in weight to \(\mu_*\) with \(D\to0\). The exterior potential \(P\to0\) by \cref{sec:exterior}. Extract one fixed clipping and pairing pattern, one chart, and a limiting core coordinate \(x_*\). By \eqref{eq:core-length}, \(E(x_*)=c\). Use the power form of the subanalytic gradient inequality \cite[Theorem~11 and Corollary~12 of the arXiv version]{BDLS} \[
 |E(x)-c|^\theta\le C\|\nabla E(x)\|
\] near \(x_*\) at the smooth observations. The Clarke subgradient there is the single classical gradient. The function shifted by its value at \(x_*\) can be restricted to a slightly larger closed bounded subanalytic chart, extended by \(+\infty\) outside, so it is lower semicontinuous and globally subanalytic, and the observations lie in the chart interior. The inequality holds for both signs of \(E-c\). At \(E=c\) it holds directly. Increasing the exponent to some rational \(0<\theta<1\) preserves the inequality near the level.

Use \(H=F-c+g_0\le |E-c|+|F-E|+g_0\). Fix a single threshold \(C_0\) large enough for the angular estimates in \cref{sec:force} when \(P>C_0d\). If \(P\le C_0d\), \eqref{eq:core-length} gives \(|E-c|\le Cd^{1/\theta}\), so \(H\le C d/\log^2(1/d)\). Otherwise \(P/d\) exceeds the fixed threshold and \eqref{eq:pairings}, \eqref{eq:core-length} give \(|E-c|^\theta\le CP\). The estimate \eqref{eq:pairings} shows \(P\le\zeta^M\) eventually for every fixed \(M\). Use Lemma~\ref{lem:separation} with \(p=P\) on the open set in the bounded chart where each endpoint pair of a geodesic piece is distinct. This set is globally subanalytic and \(E\) is smooth there. The lemma and the gradient residual bound give
\[
c_4 P\bigl(\log(1/P)\bigr)^{-(1-\varepsilon)}\le d < P/C_0.
\]
The logarithms of \(1/d\) and \(1/P\) are therefore comparable as \(P\to0\). It follows that \[
 \frac{P}{\log^2(1/P)}\le C\,\frac{d}{\log^{1+\varepsilon}(1/d)} .
\] Both the energy error in \eqref{eq:core-length} and \(|E-c|\le C P^{1/\theta}\) are bounded by the right side for small \(P,d\). Since \(B+g_0\le d^2\), both regimes yield \begin{equation}\label{eq:log-estimate}
 H\le C\,\frac{d}{\log^{1+\varepsilon}(1/d)}.
\end{equation}

This inequality holds uniformly with some fixed \(C,\varepsilon>0\) for all sufficiently late low-defect slices in a weak weight neighborhood of \(\mu_*\). Indeed, take the sequence of proposed neighborhoods given by the first \(j\) tests of a countable dense family of \(C_c^1\) tests, within \(1/j\) of their values at \(\mu_*\). Use time threshold \(j\), defect threshold \(1/j\), proposed constant \(j\), and gain \(1/j\). If all those proposals failed, contrary observations would have \(\mu_\tau\to\mu_*\) by the length bounds and \(D\to0\). The subsequence argument of this section yields some constant and a fixed positive exponent gain, contradicting the chosen failures. The same reasoning for \(P\to0\) ensures any needed initial smallness after shrinking the neighborhood and low-defect threshold.

Put \(\omega(s)=s(\log(e/s))^{1+\varepsilon}\) for small positive \(s\) and \(\omega(0)=0\). In the chosen neighborhood, \eqref{eq:log-estimate} and monotonicity of \(\omega\) near zero give \(\omega(H)\le C(\sqrt D+\eta)\). At slices with \(D\) above the fixed threshold this inequality holds for late times because \(H\to0\). For \(\Psi(s)=\int_0^s du/\omega(u)\), the integral is finite and tends to zero at zero. Where \(H>0\) and the flow is in the neighborhood, \[
 -\frac{d}{d\tau}\Psi(H)
 \ge \frac{D}{C(\sqrt D+\eta)}
 \ge C^{-1}(\sqrt D-\eta).
\] If \(H\) reaches zero, the nonnegativity and \(H'\le-D\) give \(D=0\) thereafter. Start at a sufficiently late time of the original good sequence in a smaller neighborhood whose finitely many test inequalities have positive margins in the chosen one. Up to a first exit the time integral of \(\sqrt D\) is at most \(C\Psi(H(\tau_j))+\int_{\tau_j}^\infty\eta\), tending to zero as \(j\to\infty\). The integral of \(D\) is already small on that tail. By \eqref{eq:weight-tests} the change in each defining test is smaller than its margin, ruling out first exit. This gives \(\int^\infty\sqrt D<\infty\). The test-function estimate \eqref{eq:weight-tests} then yields convergence of all the spatial weights to \(\mu_*\) in fixed coordinates.

\section{Tangents at arbitrary times}\label{sec:tangents}

\subsection{Short comparison windows}\label{sec:windows}

Take any sequence of times \(t_i\uparrow T\) whose rescaled arclength varifolds converge locally to \(V\). Write \(\tau=-\frac12\log(T-t)\), \(y=(X-x_0)/\sqrt{2(T-t)}\) and \(D(\tau)=\int (k+y\cdot n)^2e^{-|y|^2/2}\,ds_y\), with \(k\) the curvature in the \(y\) variables. The energy identity in \cref{sec:gaussian} gives \(\int^\infty D<\infty\). Set \(b_i=\int_{\tau_i}^{\tau_i+1}D\) and \(d_i=\max\{\sqrt{b_i},1/i\}\) for late \(i\). Choose \(\sigma_i\in(\tau_i,\tau_i+d_i)\) with \(D(\sigma_i)\le 2b_i/d_i\). If \(b_i=0\), use any interior time with zero defect. Then \(D(\sigma_i)\to0\).

For each \(i\) keep the physical normalization at \(t_i\) fixed throughout the window. The spatial variable is \(z=(X-x_0)/\sqrt{2(T-t_i)}\), and the time variable is \(t=(t_{\rm phys}-t_i)/(2(T-t_i))\). The window is \([0,L_i]\), where \(L_i=(1-e^{-2(\sigma_i-\tau_i)})/2\to0\). Write \(\Gamma_i(t)\) and \(\mu_i(t)\) for the graph and its arclength, and use \(k\) for curvature in this fixed spatial scale below. For each fixed compact spatial set \(U\) the following hold with constants depending on its enlargement and the original flow.

\begin{itemize}
\tightlist
\item
  There is a centered length bound \(\mu_i(t)(B_r(z))\le C r\) for all centers \(z\in U\), \(0<r\le1\), and \(0\le t\le L_i\). In physical variables the radius is \(\widehat r=r\sqrt{2(T-t_i)}\), and the centers tend to \(x_0\) uniformly on \(U\). The centered physical estimate in \cref{sec:gaussian} bounds the physical length in the ball by \(C\widehat r\). Dividing by the scale factor gives \(Cr\).
\item
  The local action \(E_i(U)=\int_0^{L_i}\int_{\Gamma_i(t)\cap U} k^2\,ds\,dt\) tends to zero, and \(\int_{\Gamma_i(L_i)\cap U}k^2\,ds\le C_U\). Indeed \(z=e^{-(\tau-\tau_i)}y\), \(dt=e^{-2(\tau-\tau_i)}d\tau\), and \(k_y^2\,ds_y\) changes by the factor \(e^{\tau-\tau_i}\) in the \(z\) variables. Use \(k_y^2\le2(k_y+y\cdot n)^2+2|y|^2\) and the compact length bound. This gives \(E_i(U)\le C_U(b_i+d_i)\) and a uniform terminal slice bound.
\item
  There are no physical endpoints in \(U\) at any time of the window for late \(i\). All degree-one endpoints are physically fixed at positive distance from \(x_0\). Every simple material cycle has length at least \(c_0>0\) in the frozen scale, uniformly throughout the windows. Each cycle encloses a bounded complementary face. A collapsing bounded face has area at least \(\pi/6\) in the instantaneous normalization by the transported area formula and quantization in multiples of \(\pi/3\) of its oriented smooth curvature. A persistent bounded face has area tending to infinity. The frozen dilation differs from the instantaneous one by a factor tending to one, so \cref{sec:faces} gives \(c_0\). A closed smooth component is included as a simple cycle. The number of true vertices and original edges is bounded by the original graph.
\end{itemize}

For a compactly supported smooth function \(\varphi\) on \(U\), balanced length variation gives \begin{equation}\label{eq:window-variation}
\frac{d}{dt}\mu_i(t)(\varphi)
 =\int_{\Gamma_i(t)} [k n\cdot\nabla\varphi-\varphi k^2]\,ds.
\end{equation} For example, integrate the tangential arclength derivative on every material edge. At a true vertex the boundary contribution is the same vertex velocity dotted with the sum of its outward unit tangents, times \(\varphi\), and vanishes. For fixed \(\varphi\), the difference between any two times in the window is \(O_\varphi(E_i(U)+\sqrt{L_i E_i(U)})\). In particular the weights \(\mu_i(t)\) converge to \(\mu=\|V\|\) uniformly over the windows on every fixed compactly supported spatial test.

In addition, the entire support avoids compact subsets of \(\mathbb R^2\setminus\operatorname{spt}\mu\) for late \(i\), uniformly over the windows. At a point of the prelimit graph in a fixed compact set, and for any sufficiently small radius \(r>0\) with the \(r\)-ball away from fixed endpoints, one can follow a nonbacktracking path of length \(r<c_0\) starting at that point. No repetition can occur before that length because it would contain a simple cycle of length at most \(r\). There is an outgoing choice at each true vertex. This path stays in the \(r\)-ball and supplies mass at least \(r\). If points at times in the windows approach a point off \(\operatorname{spt}\mu\), a fixed test supported off that support detects positive mass, contrary to uniform weight convergence.

At \(t=L_i\), the comparison window ends at the physical time corresponding to \(\tau=\sigma_i\). The two normalizations there differ by the factor \(e^{-(\sigma_i-\tau_i)}\to1\) from \(y\) to \(z\). Applying \cref{sec:compactness} at \(\sigma_i\) gives, after extraction, a Gaussian-stationary integral varifold with weight \(\mu\), since \(D(\sigma_i)\to0\) and the weights converge uniformly in the windows. At \(\mu\)-almost every point, its support is locally one smooth arc with constant positive integer weight \(m\).

\subsection{Persistence of a true junction under slice comparison}\label{sec:vertices}

\begin{lemma}\label{lem:vertices}
Fix a bound for the finite number of true vertices and edges, an upper mass bound in \(B_8(0)\), and a positive lower bound on simple cycle lengths. Consider two finite embedded network slices \(A,B\) satisfying these bounds, with \(120^\circ\) triple junctions and no physical endpoints in \(B_8(0)\). Set \(\|\varphi\|_{C^1}=\|\varphi\|_\infty+\|\nabla\varphi\|_\infty\). For every fixed \(\eta\in(0,1/2)\) there is \(\delta>0\) with the following property. If each slice has \(\int_{B_8}k^2\,ds\le\delta\) and
\begin{equation}\label{eq:slice-weight}
\sup_{\varphi\in C_c^1(B_8),\,\|\varphi\|_{C^1}\le1}
     |\mu_A(\varphi)-\mu_B(\varphi)|\le\delta,
\end{equation}
then for every true vertex \(v_A\in B_1(0)\) there is a true vertex of \(B\) within distance \(\eta\) of \(v_A\).

\end{lemma}

\begin{proof}
Suppose otherwise for a sequence with \(\delta\downarrow0\). Extract \(v_A\to p\in\overline{B_1}\) and a common weak limiting weight in \(B_8\). Slice \(B\) has no true vertices in \(B_{\eta/2}(p)\) eventually. Components in \(B_{\eta/2}(p)\) reaching \(B_{\eta/4}(p)\) have length bounded below by the collar width. A component entirely contained in this ball would have to be closed. Its tangent oscillation is bounded by \((\mu_B(B_8)\int_{B_8}k^2\,ds)^{1/2}=o(1)\), which rules out a closed curve. Thus the retained components are finite in uniformly bounded number and converge, after extraction in arclength, to straight segments across a neighborhood of any of their interior points near \(p\). This accounts for the full weight near \(p\). If two distinct limiting lines pass through \(p\) with different directions, choose a small disk around \(p\) inside \(B_{\eta/4}(p)\). The prelimit arcs then cut its boundary in alternating order, which forces a planar crossing. This is impossible for embedded vertex-free arcs. Other limiting lines have positive distance from \(p\) because there are finitely many. The common weight therefore has either empty support near \(p\) or the support of one straight line there, possibly with multiplicity.

Collect the vertices of \(A\) reachable from \(v_A\) by material edges whose lengths tend to zero. Passing to subsequences fixes the local incidence choices and makes every relevant full edge length converge in \([0,\infty]\). The collapsed cluster is a tree by the cycle length bound, and contains no physical endpoints. At the two ends of every collapsed edge the limiting outward tangents are opposite by the curvature bound. Every incident noncollapsed arm admits a fixed short initial arclength piece in the interior of \(B_8\). Its \(C^1\) limit has positive length, is supported by the common weight, and is tangent to the line at \(p\) if that line is present. If the support were empty, there would be no such arm, impossible in a finite tree of trivalent vertices. In the line case, delete each tree edge in turn and sum the limiting force balances on one side. The contributions from the interior collapsed edges cancel, so the tangent of the deleted edge is parallel to the line. This would make all three outgoing tangents at \(v_A\) parallel and contradict the true \(120^\circ\) condition. This proves the comparison statement. The comparison locates a junction of \(B\) geometrically and does not require a particular material label.
\end{proof}

The tree portion of the proof works as well for a sequence with a uniformly bounded local curvature integral whose weights converge locally to arclength supported on a smooth arc, with positive multiplicity. A collapsed edge still has tangent oscillation at most the square root of its length times the curvature integral. The noncollapsed arms have \(C^1\) limits by the local curvature bound and have tangents along the smooth arc. Therefore a sequence of such slices has no true vertices eventually on every smaller compact neighborhood meeting only that smooth arc. The empty-support variant follows by the noncollapsed arm argument. In particular this exclusion applies to the terminal slices \(\Gamma_i(L_i)\).

\subsection{An action estimate for the projections of the junctions}\label{sec:sweep}

Choose an open rectangular patch \(U\), a compact smaller rectangle \(Q\subset U\), a unit projection direction \(e\), and a Cartesian coordinate \(\xi=z\cdot e\). Assume the terminal slices have no true vertices on a fixed neighborhood of \(Q\) contained in \(U\). Define a positive measure on the projection axis by \[
\nu_i(I)=\int_0^{L_i}\int_{\Gamma_i(t)\cap U\cap\{\xi\in I\}} k^2\,ds\,dt,
\qquad
\mathcal B_i=\left\{\xi:\sup_{a>0}\frac{\nu_i([\xi-a,\xi+a])}{2a}>\epsilon\right\}.
\] Here \(\epsilon>0\) will be fixed sufficiently small. The one-dimensional covering bound gives \(|\mathcal B_i|\le C E_i(U)/\epsilon\). Slightly enlarge each witnessing interval to an open interval preserving the strict inequality, which also shows that the bad set is open. On each compact subset take finitely many open witnesses. Successively select a longest remaining interval and discard intervals meeting it. The selected intervals are disjoint and their triple dilations cover the compact subset. Use inner regularity of length. The sum of their lengths is bounded by the total measure divided by \(\epsilon\). In particular, if \(p\in Q\) has \(p\cdot e\notin\mathcal B_i\), every disk \(B_{8r}(p)\subset U\) satisfies \(E_i(B_{8r}(p))\le16\epsilon r\).

For late \(i\), every true vertex \(p\in Q\) at any time \(t_0\in[0,L_i]\) has \(p\cdot e\in\mathcal B_i\). Suppose instead that it projects outside \(\mathcal B_i\). Put \(r_0=\sqrt{L_i/2}\) and \(r_j=2^{-j}r_0\). We assume \(i\) large enough that \(B_{8r_0}(p)\subset U\) and \(B_{r_0}(p)\) lies in the terminal vertex-free neighborhood of \(Q\), uniformly for \(p\in Q\). Set \(s_0=L_i\). For each \(j\ge1\), there is an interval of length \(r_j^2\) starting or ending at \(t_0\) inside the window, since one time direction has length at least \(L_i/2\). Choose a time \(s_j\) inside that interval with \[
r_j\int_{\Gamma_i(s_j)\cap B_{8r_j}(p)} k^2\,ds\le 32\epsilon .
\] This is the curvature-action average on that interval. At a sufficiently large final index \(J\), take \(s_J=t_0\) itself. At the fixed pre-singular slice \(\Gamma_i(t_0)\) the total curvature integral is finite, so its product with \(r_J\) can be made at most \(32\epsilon\). We need only the resulting finite list of times, which need not be ordered. At \(s_0\) the local curvature integral is bounded uniformly in \(i\), so its product with \(r_0\) tends to zero.

Compare the slices at \(s_j\) and \(s_{j-1}\) in the spatial scale \(r_j\) centered at \(p\). They have small curvature integrals in \(B_8\) there, by containment in the coarser averaging disk and \(r_j=r_{j-1}/2\). The rescaled masses in \(B_8\) have a common upper bound from the local length ratios. The cycle length lower bound remains valid in these magnifications and all fixed endpoints stay outside \(B_8\). For \(\varphi\) as in \eqref{eq:slice-weight}, set \(\varphi_j(z)=\varphi((z-p)/r_j)\). The magnified weight at time \(s\) evaluates \(\varphi\) as \(r_j^{-1}\mu_i(s)(\varphi_j)\). With \(\Delta\mu_i=\mu_i(s_j)-\mu_i(s_{j-1})\), the variation formula \eqref{eq:window-variation} gives \begin{equation}\label{eq:comparison}
\begin{split}
r_j^{-1}\left|\Delta\mu_i(\varphi_j)\right|
&\le r_j^{-2}\big(Cr_j|s_j-s_{j-1}|\,E_i(B_{8r_j}(p))\big)^{1/2}
        +r_j^{-1}E_i(B_{8r_j}(p))\\
&\le C(\sqrt\epsilon+\epsilon).
\end{split}
\end{equation} The time separation is at most \(9r_j^2\), also for \(j=1\) because \(|s_0-t_0|\le L_i=2r_0^2\). The estimate \eqref{eq:comparison} uses curvature action on the whole window and applies in either time order.

Take \(\eta=1/4\) in Lemma~\ref{lem:vertices} and choose \(\epsilon\) so both smallness conditions hold. At \(s_J\) there is a true vertex at \(p\). If at \(s_j\) there is a vertex in \(B_{r_j}(p)\), \eqref{eq:slice-weight}--\eqref{eq:comparison} give one at \(s_{j-1}\) in \(B_{(1+\eta)r_j}(p)\subset B_{r_{j-1}}(p)\). Repeating the finite induction gives a terminal vertex in \(B_{r_0}(p)\), a contradiction. Thus \begin{equation}\label{eq:sweep}
\left|\operatorname{proj}_{e}
 \{p\in Q:p\text{ is a true vertex at some }t\in[0,L_i]\}\right|
 \le C\epsilon^{-1}E_i(U)\longrightarrow0.
\end{equation} The scale-dependent curvature comparisons in this proof involve every true vertex that projects outside the bad intervals. A track may move quickly, and the argument never estimates its speed or assumes that one material vertex persists between chosen comparison slices.

\subsection{Projected intersection counts}\label{sec:counts}

Fix an isolated smooth arc of \(\operatorname{spt}\mu\) on which its weight is \(m\ge1\). For any direction \(e\) with nonzero tangential component at the chosen point, take a small rectangular strip \(Q=\{\xi\in[a,b],\,h_-\le z\cdot e^\perp\le h_+\}\) around the point. In a slightly larger patch \(U\), the arc is a graph over \(\xi\), crosses the whole projection interval, and is separated from the horizontal boundaries of the strip. There is no other limit support in the strip or at its horizontal boundaries. With these strict margins, uniform support convergence in \cref{sec:windows} excludes all crossings of the horizontal boundaries throughout the windows. The terminal slice exclusion in \cref{sec:vertices} forbids true vertices at \(t=L_i\) on a neighborhood of \(Q\).

The terminal tangents converge uniformly in the strip, as unoriented lines, to the tangent of the local limiting graph at the corresponding projected coordinate. Otherwise one can choose a contrary sequence of centered arclength segments with fixed positive length on either side, contained in the larger vertex-free patch. Their \(C^1\) limit from the local curvature bound would be supported by the single smooth arc but have the wrong tangent. The projection derivative on every terminal strand is consequently bounded away from zero in absolute value. Every component meeting the interior of the strip crosses from \(\xi=a\) to \(\xi=b\) because it has no interior endpoint or vertex, no horizontal escape, and strictly monotone projection. There are finitely many such through-arcs by the length bound. Let \(M_i\) be their common crossing count. Take a nonzero nonnegative smooth test \(\psi\) supported in \((a,b)\) and a spatial cutoff equal to one near the limiting arc over that support and supported inside the band. Uniform support avoidance puts all the terminal strands over \(\operatorname{spt}\psi\) in the part where the cutoff is one. The projection formula proved in \cref{sec:distribution} gives
\[
M_i\int_a^b\psi(x)\,dx
 =\int_{\Gamma_i(L_i)\cap Q}\psi(z\cdot e)|u_i\cdot e|\,ds,
\]
where \(u_i\) is the terminal unit tangent and the integrand includes the cutoff. Uniform tangent convergence and weak weight convergence give \(m\int_a^b\psi(x)\,dx\) in the limit by the same formula on the graph. Hence \(M_i\to m\), and \(M_i=m\) for late \(i\) since both are integers.

Fix a projected coordinate \(x\in(a,b)\setminus\mathcal B_i\). The covering estimate in \cref{sec:sweep} excludes true vertices on \(\xi=x\) inside the strip throughout the window. The horizontal boundaries and genuine fixed endpoints are also disjoint from its intersection with the network. At almost every such \(x\) the starting slice is transverse. A \(C^1\) projection of a compact material parameter interval maps its critical set to a null set: cover the critical set by the relatively open set where the derivative magnitude is below \(\epsilon'\). The interval components have total length at most the parameter length and the sum of their image lengths is at most \(\epsilon'\) times that length. Let \(\epsilon'\) tend to zero and exclude the finitely many endpoint images. The number of intersections of \(\xi=x\) in the strip cannot increase from that starting value. In fact, in any smooth fixed material parametrization \(X(u,t)\) of an edge, the function \(g(u,t)=X(u,t)\cdot e-x\) solves \[
g_t=|X_u|^{-2}g_{uu}
   +\left(\lambda_{\rm tan}/|X_u|-\frac{X_u\cdot X_{uu}}{|X_u|^4}\right)g_u,
\] where \(\lambda_{\rm tan}\) is the tangential velocity. On compact interior parameter-time rectangles this is a smooth uniformly parabolic scalar equation. Angenent's zero-number theorem \cite[Theorem~A]{AngNodal} applies with nonzero lateral boundary values on such rectangles. Shift time to the beginning of each rectangle. The coefficients \(a,b,c\) are respectively \(|X_u|^{-2}\), the displayed coefficient of \(g_u\), and zero. Smoothness and \(|X_u|>0\) on each compact rectangle ensure all its coefficient derivatives and strict positivity.

Here is the localization at an arbitrary interior time \(t_*\). The set of parameter zeros of \(g\) in the band at \(t_*\) is compact and avoids the parameter endpoints and the horizontal boundaries. Each connected component of the full zero set that meets the band is wholly in its interior, since there is no boundary intersection. It is an interval or a point with nonzero values of \(g\) arbitrarily close on either side. An entire closed component on the cutting line inside the strip is impossible by embeddedness. Thus the in-band zero set is covered by a finite union of interior parameter intervals whose endpoints have nonzero \(g\) and whose closures map inside the open band at \(t_*\). Merge overlapping intervals, retaining nonzero outer endpoints. By compactness and continuity, on a time neighborhood of \(t_*\) those intervals capture all in-band zeros, have nonzero boundary values, and have no out-of-band zeros. Apply the zero-number theorem to their sum of distinct zeros. This gives local nonincrease in time, including degenerate folds. At time zero and the terminal time the roots are transverse and no endpoint crossing is possible, giving the same comparison up to the endpoints. The positive-time geometric representatives described in \cref{sec:introduction} supply these local regular parametrizations. Periodic edge parameters can be cut away from the in-band zero set. The material parametrizations need only be smooth within each fixed comparison window and on the indicated interior rectangles. No uniform lower bound on edge lengths or uniform ellipticity as \(i\to\infty\) is used.

It follows that \(N_i(0,x)\ge N_i(L_i,x)=m\) for almost every \(x\) outside \(\mathcal B_i\), where \(N_i\) counts intersections of the vertical cut in the strip with arclength multiplicities. In particular the measure of levels with fewer than \(m\) intersections satisfies \begin{equation}\label{eq:deficit}
\big|\{x\in(a,b): N_i(0,x)<m\}\big|\le C\epsilon^{-1}E_i(U).
\end{equation} The right side tends to zero for each fixed strip. It bounds the lost set of projected coordinates, allowing for folds of any strand and for changes in the participating arms at true junctions.

\subsection{The limiting tangent distribution}\label{sec:distribution}

The one-dimensional projection formula used here and in \cref{sec:counts} follows by splitting the open set where the projection derivative is nonzero into disjoint monotonicity intervals and changing variables on each. The length integrand weighted by the absolute derivative vanishes on the critical set, whose image is null by \cref{sec:counts}. Let \(\psi\ge0\) be a smooth test compactly supported in \((a,b)\). With the spatial band cutoff described above, the projection formula and \eqref{eq:deficit} give \begin{equation}\label{eq:tangent-order}
\int \psi(z\cdot e)|u\cdot e|\,dV(z,[u])
\ \ge\ m\int_a^b\psi(x)\,dx
\ =\ \int \psi(z\cdot e)|u_*(z)\cdot e|\,d\mu(z),
\end{equation} where both spatial integrands include the band cutoff and \(u_*\) is a unit tangent of the limit arc. To justify the first step, before passage to the limit the first integral is \(\int\psi(x)N_i(0,x)\,dx\). The lower bound on that integral loses at most \(m\|\psi\|_\infty|\mathcal B_i|\). The angular integrand \(|u\cdot e|\) is continuous on unoriented lines. The right equality is coarea on the limit arc.

On a patch where \(z\mapsto z\cdot e\) is one-to-one on the regular limit arc, \eqref{eq:tangent-order} yields \[
\int_{\{z\in A\}}|u\cdot e|\,dV\ \ge\ \int_A|u_*(z)\cdot e|\,d\mu
\] for every Borel subset \(A\) of that arc patch. Indeed the spatial measures are supported on the arc, the projection is a homeomorphism onto its interval, and nonnegative smooth tests determine measure order there. For each \(e\) in a countable dense set of directions, take countably many such patches covering the regular points where \(|u_*\cdot e|>0\) and restrict their inequalities to disjoint Borel pieces. At the complementary regular points with \(|u_*\cdot e|=0\) the inequality is automatic. The exceptional points outside the regular locus have zero \(\mu\)-measure and hence zero \(V\)-measure over them.

For each \(\epsilon'>0\), partition the regular locus into Borel sets \(A_j\) where the unit tangent is within \(\epsilon'\), up to sign, of the corresponding direction \(e_j\) from the dense set, using the least such index. On \(A_j\), one has \(||u\cdot u_*|-|u\cdot e_j||\le\epsilon'\) and \(|u_*\cdot e_j|\ge1-\epsilon'\). For any spatial compact \(K\) we conclude \[
\int_{\{z\in K\}}(1-|u\cdot u_*(z)|)\,dV\le 2\epsilon'\sum_j\mu(K\cap A_j)=2\epsilon'\mu(K).
\] Letting \(\epsilon'\downarrow0\) identifies \([u]\) with the tangent line of \(\mu\) almost everywhere in \(V\). If \(\mu=0\), then \(V=0\).

\begin{proof}[Completion of the proof of \cref{thm:main}]
The length bounds in \cref{sec:gaussian} and compactness of \(G(2,1)\) give local weak precompactness of the rescaled varifolds. The weight of any subsequential limit at arbitrary times is \(\mu_*\) by \cref{sec:trapping}. The tangent identification in \cref{sec:distribution} puts its full mass on the tangent line of this weight, with the integer multiplicity obtained in \cref{sec:compactness}. Thus \(V=V_*\) as measures on position and unoriented tangent line. The scale and the coordinates in \eqref{eq:rescale} are fixed for every subsequence.
\end{proof}


\section*{AI assistance}
AI tools assisted with proof development, source checking, and
manuscript preparation. Separate AI roles assessed the complete
natural-language argument and reviewed a manuscript draft. They
worked in a shared environment with some prior task context. The
draft was revised after the reader and referee reports.

\begin{thebibliography}{99}

\bibitem{AIM}
S. Brunton, M. J. Colbrook, M. V. de Hoop, G. Stepaniants,
A. Townsend, and R. Ward,
\emph{159. Uniqueness of the limiting shape at a network singularity},
in \emph{AIM --- Open Applied Problems}, 2026.
\href{https://github.com/MColbrook/AIM/blob/b2ec0fac559f532fc0ed04fce81deea0ce07c68b/problems/159-network-tangent-uniqueness.md}{Fixed problem statement}.


\bibitem{AngCurve}
S. B. Angenent,
\emph{Parabolic equations for curves on surfaces. Part I. Curves with
\(p\)-integrable curvature},
Ann. of Math. (2) \textbf{132} (1990), 451--483.
\href{https://doi.org/10.2307/1971426}{doi:10.2307/1971426}.

\bibitem{AngNodal}
S. B. Angenent,
\emph{Nodal properties of solutions of parabolic equations},
Rocky Mountain J. Math. \textbf{21} (1991), 585--592.
\href{https://doi.org/10.1216/rmjm/1181072953}{doi:10.1216/rmjm/1181072953}.

\bibitem{BDLS}
J. Bolte, A. Daniilidis, A. S. Lewis, and M. Shiota,
\emph{Clarke subgradients of stratifiable functions},
SIAM J. Optim. \textbf{18} (2007), 556--572.
\href{https://doi.org/10.1137/060670080}{doi:10.1137/060670080}.
\href{https://arxiv.org/abs/math/0601530v1}{arXiv:math/0601530v1}.

\bibitem{CL}
J.-E. Chang and Y.-K. Lue,
\emph{Uniqueness of regular shrinkers with 2 closed regions},
\href{https://arxiv.org/abs/1901.10315v1}{arXiv:1901.10315v1} (2019).

\bibitem{MNPS}
C. Mantegazza, M. Novaga, A. Pluda, and F. Schulze,
\emph{Evolution of networks with multiple junctions},
Astérisque \textbf{452} (2024).
\href{https://people.dm.unipi.it/novaga/papers/mantegazza/tripunto/network-survey.pdf}{Author version linked by AIM}.
\href{https://arxiv.org/abs/1611.08254v4}{arXiv:1611.08254v4}.

\bibitem{PP}
A. Pluda and M. Pozzetta,
\emph{On the uniqueness of nondegenerate blowups for the motion by curvature of networks},
\href{https://arxiv.org/abs/2204.11675v2}{arXiv:2204.11675v2} (2022).

\bibitem{Teschl}
G. Teschl,
\emph{Ordinary Differential Equations and Dynamical Systems},
Graduate Studies in Mathematics, vol. 140, Amer. Math. Soc.,
Providence, RI, 2012.
\href{https://mat.univie.ac.at/~gerald/ftp/book-ode/ode.pdf}{Author's version}.

\bibitem{Valette}
G. Valette,
\emph{On subanalytic geometry},
\href{https://arxiv.org/abs/2507.23622v1}{arXiv:2507.23622v1} (2025).

\end{thebibliography}
\end{document}
